\scanpage{57}
\mspage{56}

\uline{Proposition} a) Si $F$ \deleted{\unsure{est|était}} est une figure, alors les figures $F'$ telles que $F'\leq F$ s'identifient, par l'application $F'\mapsto\overline F'$, aux parties de $\overline F$ qui sont \uline{fermées} pour la relation d'ordre \inserted{$\leq$} de $\unsure{\mathcal M|\mathscr M}$.

b) Si $F$ et $G$ sont deux figures, $F\wedge G$ existe (inf pour $\leq$) \inserted{\uuline{figure intersection}} et \deleted{\unsure{on|en}} $\overline{F\cap G}=\overline F\cap\overline G$. Même pour une famille $(F_i)_{i\in I}$ \unsure{quelc.|qelc.} \inserted{avec \uuline{$I\neq\varnothing$}}.

% margin note written sideways
\margin{\uline{Mais attention} il vaut mieux ne pas les noter $F\cap G$, car quand il y a une notion de support, ça ne correspond pas à l'intersection des supports, sauf si $F,G$ compatibles.}

[ \uline{Lieux} de $\unsure{\mathcal F|\mathscr F}$ : ce sont les \deleted{\unsure{pas|pet}} figures qui sont \uline{minimales} dans $\mathcal F\setminus\{\varnothing_{\mathcal F}\}$. \deleted{\unsure{Un|Une},} \deleted{Un lieux est de \unsure{ce nom|nom}}  Tout lieu est irréductible. Donc les lieux sont aussi les éléments minimaux de $\mathcal M$ : ce sont les multistrates qui n'ont d'autres strates \unsure{qu'elles-même|qu'elles-mêmes|qu'elle-même}. ] \deleted{Si $F$ est une figure, les}

% slanted
\margin{Définition \unsure{courante|courte} n'est raisonnable que dans certains cas \deleted{\unsure{(quand les|(quand la}} (p. ex. les contrées \unsure{en boules''|en boules}, ou contrées \unsure{par morceaux|ponctuées} $\ldots$) \unsure{Car dans|Car chez} les \unsure{cas|cat.} types d'\unsure{ordre|données} \unsure{quelconq.|quelconques}, \unsure{un|une} \unsure{modèle|modélé} $\overline E$ \unsure{de str.|de structure} \unsure{semi-}\illegible{1} \inserted{\unsure{univers vides|univers vide},} forme \deleted{\unsure{peut être une|peut être un}} \inserted{une} multistrate \unsure{minimale|minime}, mais on ne \unsure{l'appelle|appelle} un ``lieu'' que si elle est \unsure{réduite|discrète|vide réelle} !}

\uline{Figures} \uline{disjointes} \unsure{,,|''} \uline{$F$} et \uline{$G$} : telles que a) $F$ et $G$ compatibles, i.e $F\cup G$ existe \deleted{et b)} ou encore $\{F,G\}$ majoré pour $\leq$ et b) $F\cap G=\varnothing_{\mathcal F}$ \deleted{\unsure{ou|et} Cela signifie \unsure{donc|que}} \inserted{(i.e. $\overline F\cap\overline G=\varnothing$)}.

% slanted
\margin{\unsure{bloque|bloquer} mes ``\unsure{coproduits|composantes}''}

\deleted{Si $F$ et $G$ sont} \uline{Prop} Si $F$ et $G$ sont disjointes, alors pour $F'\leq F$, $G'\leq G$, $\unsure{F'|F'_{a}}$ et $G'$ sont disjointes, et $F'\cup G'$ est une sous figure de $F\cup G$, l'application $(F';G')\mapsto F'\cup G'$ est une bijection \deleted{de l'\unsure{ens|ens.} de $\unsure{\mathcal F|F}$}
